% Allow the overview to float near the section opening without blocking prose.
% Shared English visual; both prose editions include this source.
\par\medskip
\noindent\begin{minipage}{\linewidth}
\captionsetup{type=figure}
\centering
\begin{tikzpicture}[
  summary/.style={draw=rulegray,rounded corners=2pt,fill=boxgray!45,
    text width=3.75cm,minimum height=4.3cm,inner sep=5pt,
    align=center,font=\small},
  heading/.style={font=\sffamily\small\bfseries,anchor=south}]
\node[summary] (source) at (0,0) {
  $\displaystyle t=\sum_j x_j,\quad \mu=t/N$\\[5pt]
  $\displaystyle ss=\sum_j(x_j-\mu)^2$\\[9pt]
  Full row: $ss=\orig{\exTwoPass}$.\\[4pt]
  Adding local results:\\[2pt]
  $ss_\pa+ss_\pb=\stale{\exNaive}$
};
\node[summary] (reference) at (4.55,0) {
  Reference state tuple: $(t,q,n)$\\[5pt]
  $\begin{aligned}
    t&=\textstyle\sum_j x_j\\[2pt]
    q&=\textstyle\sum_j x_j^2\\[2pt]
    n&=\textstyle\sum_j 1
  \end{aligned}$\\[5pt]
  Componentwise join\\[3pt]
  $ss=q-t^2/n$\\[5pt]
  Shifted float32: $\stale{\exTextbookFloat}$
};
\node[summary] (maintained) at (9.1,0) {
  Maintained state tuple: $(t,ss,\aux{n})$\\[5pt]
  $\mu=t/n,\quad \mu_p=t_p/n_p$\\[5pt]
  $ss=ss_\pa+ss_\pb+\Delta$\\[3pt]
  $\displaystyle\Delta=\sum_{p\in\{\pa,\pb\}}n_p(\mu-\mu_p)^2$\\[7pt]
  Shifted float32: $\fix{\exTwoPassFloat}$
};
\node[heading,text=origblue] at ([yshift=3pt]source.north) {(a) Two-pass source};
\node[heading] at ([yshift=3pt]reference.north) {(b) Reference states};
\node[heading,text=fixgreen] at ([yshift=3pt]maintained.north) {(c) Maintained states};
\draw[flow] (source.east) -- (reference.west);
\draw[flow,draw=fixgreen] (reference.east) -- (maintained.west);
\end{tikzpicture}
\caption{Deriving a single-pass variance reducer.
For $\pa=(1,3)$ and $\pb=(2,6)$, adding locally centred results loses the change of mean.
The reference state tuple establishes a three-state bound; the maintained state tuple retains $t$ and $ss$ and adds \aux{$n$} for the correction.
The shifted results report $ss$ after adding $10^4$ to every input.}
\label{fig:variance}
% src: data/illus.json variance_exact and variance_float32
\end{minipage}\par\medskip

\begingroup
\setlength{\intextsep}{2pt}
\reservewrapspace{12\baselineskip}
\section{Motivating Example}
\label{sec:example}

\begin{wrapfigure}{r}{0.48\textwidth}
\begin{lstlisting}[style=pseudocode,caption={Two-pass source program.},label={lst:variance-before}]
Input: $x_0,\ldots,x_{N-1}$, $N>0$
$t\gets0$
for $j=0,\ldots,N-1$:
    $t\gets t+x_j$
$\mu\gets t/N$
$ss\gets0$
for $j=0,\ldots,N-1$:
    $ss\gets ss+(x_j-\mu)^2$
return $\;ss/N$
\end{lstlisting}
\end{wrapfigure}
Variance illustrates the dependency that fusion must remove.
For a row of length $N$, the source program in \cref{lst:variance-before} first computes the total $t$ and the mean $\mu=t/N$, then accumulates the centred sum of squares $ss$; the variance is $ss/N$.
The second traversal cannot begin its original computation until the first has produced the final mean.

To replace this dependency with a join, we first consider results computed using local means.
For a non-empty part $S$, let $t_S$ and $ss_S$ be its source reduction results and $\mu_S$ its mean.
Computing these results directly still takes two passes over that part.
For disjoint parts $\pa$ and $\pb$, applying the source reduction operator, $\oplus=+$ in this example, gives $ss_\pa\oplus ss_\pb=ss_\pa+ss_\pb$.
This does not generally equal the centred sum for their union: the two local results use different means, as \cref{fig:variance}(a) illustrates.
A correct join must account for the change to a common mean.

We therefore first seek states that can be combined without sharing a mean.
Expanding the centred reduction makes its dependence on $\mu_S$ explicit:
\begin{equation}
\label{eq:variance-reference}
ss_S=\sum_{j\in S}(x_j-\mu_S)^2
=\underbracket{\colorbox{boxgray}{$\displaystyle\sum_{j\in S}x_j^2$}}_{q_S}
 -2\mu_S\underbracket{\colorbox{boxgray}{$\displaystyle\sum_{j\in S}x_j$}}_{t_S}
 +\mu_S^2\underbracket{\colorbox{boxgray}{$\displaystyle\sum_{j\in S}1$}}_{n_S}.
\end{equation}
The highlighted reductions do not read $\mu_S$: they are the square sum $q_S$, the total $t_S$ and the count $n_S$.
They form the reference state tuple $(t_S,q_S,n_S)$ in panel (b).
Each state combines by addition, and together they recover the source result through $\mu_S=t_S/n_S$ and $ss_S=q_S-t_S^2/n_S$.
This establishes a sufficient three-state representation, independent of the input length.

The reference representation makes a single pass algebraically possible, but its reconstruction of $ss$ subtracts potentially large, nearly equal quantities.
The shifted float32 input in \cref{fig:variance} illustrates the resulting cancellation.
To retain the centred computation, \sys instead starts from the source results $(t,ss)$ and analyses what their correction requires.
Let $\mu$ be the mean of the union and $\delta_\pa=\mu-\mu_\pa$ its change on $\pa$.
Expanding around the local mean gives
\begin{equation}
\label{eq:variance-rebase}
\sum_{j\in\pa}(x_j-\mu)^2
=ss_\pa-2\delta_\pa\sum_{j\in\pa}(x_j-\mu_\pa)+\delta_\pa^2 n_\pa
=ss_\pa+\delta_\pa^2 n_\pa.
\end{equation}
The linear term vanishes because deviations from a partition's own mean sum to zero.
The remaining correction is weighted by $n_\pa$, the number of elements whose centre changes.
The source results $(t_\pa,ss_\pa)$ do not retain this count.
Thus context-correction analysis exposes $n$ as information missing from the chosen representation, and \sys tests it as an auxiliary state.
Here that coefficient alone supplies a sufficient augmentation: since $n_S>0$, the maintained state tuple $(t_S,ss_S,n_S)$ recovers $q_S=ss_S+t_S^2/n_S$, with the other reference states read directly.
The state-selection procedure verifies this recovery before accepting the tuple.

\reservewrapspace{14\baselineskip}
\begin{wrapfigure}{r}{0.48\textwidth}
\begin{lstlisting}[style=pseudocode,caption={Fused single-pass program.},label={lst:variance-after}]
$\leaf(x)=(x,0,1)$
function $\Mrg(\pa,\pb)$:
    $t\gets t_\pa+t_\pb$; $n\gets n_\pa+n_\pb$
    $\mu\gets t/n$; $ss\gets ss_\pa+ss_\pb$
    for $p\in\{\pa,\pb\}$:
        $ss\gets ss+n_p(\mu-t_p/n_p)^2$
    return $\;(t,ss,n)$
$(t,ss,n)\gets\leaf(x_0)$
for $j=1,\ldots,N-1$:
    $(t,ss,n)\gets\Mrg((t,ss,n),\leaf(x_j))$
return $\;ss/n$
\end{lstlisting}
\end{wrapfigure}
Using the recovery to derive a join and expressing it around the local means gives panel (c).
The join adds totals and counts, and corrects each local centred sum to the combined mean before adding them.
\sys verifies that the result is the maintained state tuple of the union.
On the illustrated shifted input, the corrected join matches the two-pass result.

To turn this join into fusion, choose each part to contain exactly one element.
Its mean is $x_j$, its centred sum is zero and its count is one, so its maintained state tuple is immediately $\leaf(x_j)=(x_j,0,1)$.
\Cref{lst:variance-after} initialises the running tuple from the first element and computes and joins each subsequent tuple as its element is read.
This removes the local traversals as well as the global dependency: the input is read once, and the final variance is $ss/n$.
Subscripts in the listing identify fields of the two input tuples.

Over the reals, the resulting join recovers the pairwise variance formula of \citet{chan} and, applied one element at a time, is algebraically equivalent to Welford's update~\citep{welford}.
The same union property also permits a parallel reduction tree.
\endgroup
